\section{Conclusion}

\label{sec:conclusion}

We have presented Hanzo ACI, a blockchain architecture purpose-built for decentralized AI compute markets. The Proof of AI (PoAI) mechanism achieves 99.7\% malicious provider detection at 0.3\% verification overhead by combining TEE attestations with statistical sampling. The dual-track consensus architecture achieves 12,000 job matches per second with 200ms latency while maintaining 2-second settlement finality. The AI-native virtual machine reduces gas costs by 73\% compared to equivalent Ethereum L2 implementations. Formal security proofs establish that honest computation is the Nash equilibrium of the PoAI game, and testnet evaluation with 100 validators and 500 workers validates the system's performance and security properties under realistic conditions.

\paragraph{Hanzo AI Chain (AIVM).}
The compute markets, TEE attestation primitives, and PoAI verification described above are deployed on the Hanzo AI Chain, the AI Virtual Machine (AIVM) that re-frames Lux's unified A-Chain (LP-134) as the open AI protocol for the Hanzo Network. AI Chain mints the native \textbf{AI} coin, hosts the on-chain HMM and agent registries, and inherits Lux Quasar 3.0 three-lane post-quantum finality (LP-105) along with QuasarGPU execution (LP-132) and QuasarSTM lane scheduling (LP-010). See the companion paper \emph{Hanzo AI Chain (AIVM)} (\texttt{\textasciitilde/work/hanzo/papers/hanzo-ai-chain/}) for the full architecture, AI tokenomics, and the 25 December 2025 activation roadmap.

\bibliographystyle{plain}
\begin{thebibliography}{30}

\bibitem{akash2020}
G.~Loli and A.~Vilenski.
\newblock Akash network: Decentralized cloud infrastructure marketplace.
\newblock \emph{Akash Whitepaper}, 2020.

\bibitem{bittensor2023}
Bittensor.
\newblock Bittensor: A peer-to-peer intelligence market.
\newblock \emph{Bittensor Whitepaper}, 2023.

\bibitem{costan2016intel}
V.~Costan and S.~Devadas.
\newblock Intel {SGX} explained.
\newblock \emph{IACR Cryptology ePrint Archive}, 2016.

\bibitem{danezis2022narwhal}
G.~Danezis, L.~Kokoris-Kogias, A.~Sonnino, and A.~Spiegelman.
\newblock Narwhal and {Tusk}: A {DAG}-based mempool and efficient {BFT} consensus.
\newblock In \emph{EuroSys}, 2022.

\bibitem{gensyn2023}
B.~Fielding and H.~de~Valence.
\newblock Gensyn: A protocol for training machine learning models.
\newblock \emph{Gensyn Whitepaper}, 2023.

\bibitem{golem2016}
J.~Petkanics.
\newblock The golem project: A decentralized computational network.
\newblock \emph{Golem Whitepaper}, 2016.

\bibitem{ionet2024}
A.~Thomas and T.~Thayyil.
\newblock io.net: The internet of {GPU}s.
\newblock \emph{io.net Whitepaper}, 2024.

\bibitem{mckinseyai2024}
McKinsey.
\newblock The state of {AI} in 2024.
\newblock \emph{McKinsey Global Survey}, 2024.

\bibitem{modulus2024}
Modulus.
\newblock Verifiable {AI} inference on-chain.
\newblock \emph{Modulus Documentation}, 2024.

\bibitem{rendernetwork2021}
J.~Urbach.
\newblock Render network: Distributed {GPU} rendering.
\newblock \emph{Render Whitepaper}, 2021.

\bibitem{risczero2023}
RISC~Zero.
\newblock {RISC Zero}: General-purpose verifiable computing.
\newblock \emph{RISC Zero Documentation}, 2023.

\bibitem{ritual2024}
Ritual.
\newblock Ritual: {AI} coprocessor for blockchains.
\newblock \emph{Ritual Documentation}, 2024.

\bibitem{sui2023}
Sui.
\newblock The {Sui} smart contracts platform.
\newblock \emph{Sui Documentation}, 2023.

\bibitem{teutsch2019truebit}
J.~Teutsch and C.~Reitwie{\ss}ner.
\newblock A scalable verification solution for blockchains.
\newblock \emph{TrueBit Whitepaper}, 2019.

\bibitem{yin2019hotstuff}
M.~Yin, D.~Malkhi, M.~K. Reiter, G.~G. Gueta, and I.~Abraham.
\newblock {HotStuff}: {BFT} consensus with linearity and responsiveness.
\newblock In \emph{PODC}, 2019.

\bibitem{zkevm2023}
Polygon.
\newblock Polygon {zkEVM}: Scaling {Ethereum} with zero-knowledge proofs.
\newblock \emph{Polygon Whitepaper}, 2023.

\bibitem{castro1999practical}
M.~Castro and B.~Liskov.
\newblock Practical {Byzantine} fault tolerance.
\newblock In \emph{OSDI}, 1999.

\bibitem{gilad2017algorand}
Y.~Gilad, R.~Hemo, S.~Micali, G.~Vlachos, and N.~Zeldovich.
\newblock Algorand: Scaling {Byzantine} agreements for cryptocurrencies.
\newblock In \emph{SOSP}, 2017.

\bibitem{wood2014ethereum}
G.~Wood.
\newblock Ethereum: A secure decentralised generalised transaction ledger.
\newblock \emph{Ethereum Yellow Paper}, 2014.

\bibitem{yakovenko2018solana}
A.~Yakovenko.
\newblock Solana: A new architecture for a high performance blockchain.
\newblock \emph{Solana Whitepaper}, 2018.

\bibitem{kwon2023vllm}
W.~Kwon, Z.~Li, S.~Zhuang, et~al.
\newblock Efficient memory management for large language model serving with {PagedAttention}.
\newblock In \emph{SOSP}, 2023.

\bibitem{bonneau2015sok}
J.~Bonneau, A.~Miller, J.~Clark, A.~Narayanan, J.~A. Kroll, and E.~W. Felten.
\newblock {SoK}: Research perspectives and challenges for {Bitcoin} and cryptocurrencies.
\newblock In \emph{IEEE S\&P}, 2015.

\bibitem{buterin2014next}
V.~Buterin.
\newblock A next-generation smart contract and decentralized application platform.
\newblock \emph{Ethereum Whitepaper}, 2014.

\bibitem{gavin2023eigenlayer}
S.~Gavin, S.~Kannan, and R.~Resnick.
\newblock {EigenLayer}: The restaking collective.
\newblock \emph{EigenLayer Whitepaper}, 2023.

\bibitem{tramer2019slalom}
F.~Tram\`er and D.~Boneh.
\newblock Slalom: Fast, verifiable and private execution of neural networks in trusted hardware.
\newblock In \emph{ICLR}, 2019.

\end{thebibliography}

